\begin{papercard}{niksirat2024}
\cf{Stages} \stagetag{H} (inferences, privacy risks and user attitudes around wearable activity data), \stagetag{S} (definition of the object of study), \stagetag{R} (how a systematic literature review is documented). A 40-page \emph{ACM Computing Surveys} article (56(7), 2024) that systematically reviews 236 studies on the utility, privacy and security of wearable activity trackers (WATs). It is an evidence base for harm mechanisms and user perceptions in a neighbouring domain (fitness tracking), and a worked example of a documented review method. Findings are about wearables and their companion apps, not about mental-health apps.

\cf{Aim and novelty} Existing surveys cover wearables in general (eye-worn glasses, EEG devices) or a single aspect; this one covers WATs only and integrates utility, privacy and security, because they are ``inherently intertwined'' and because older surveys miss the last five years \pc{niksirat2024}{3}{the holistic approach also allowed us to draw a research agenda for WATs by identifying six research gaps}. Three contributions: main privacy issues from WAT use, utility aspects with privacy implications, and security vulnerabilities/attacks/countermeasures.

\cf{Review method} A systematic literature review following Kitchenham et al.'s assessment criteria \pc{niksirat2024}{3}{We conducted a systematic literature review following the assessment criteria of Kitchenham et al.}. Search string = (physical activity data OR fitness data OR fitness tracking OR wearable activity tracking OR physical activity tracker OR fitness tracker OR wearable activity tracker) AND (utility OR privacy OR security OR perception OR understanding OR experience OR expectation OR sharing) AND (system OR device OR application OR app OR service OR bracelet OR wrist-worn); keywords chosen from earlier surveys. Sources: ACM DL, IEEE Xplore, AIS library, USENIX, PoPETs, ScienceDirect, Springer Link, Google Scholar for other publishers, plus recent proceedings tracked during the review. 688 papers after deduplication; excluded non-English, published 2012 or earlier, not peer reviewed (position papers, letters, editorials, theses, patents, books) \pc{niksirat2024}{4}{(iii) are not peer reviewed (e.g., position papers, letters, editorials, prefaces, article summaries, theses, patents, or books)}; included only papers about WATs (or with implications for them) and directly relevant to utility, privacy and/or security. ``The first author applied the inclusion and exclusion criteria and later confirmed them with the second and third authors''; \textbf{236 of 688} were selected \pc{niksirat2024}{4}{The first author applied the inclusion and exclusion criteria and later confirmed them with the second and third authors}. Three authors summarised context, methodology and findings, discussed weekly; synthesis followed the JBI Manual for Evidence Synthesis (identify patterns and homogeneity, then heterogeneity) \pc{niksirat2024}{4}{we adhered to the guidelines provided in the JBI Manual for Evidence Synthesis}. The list of reviewed papers and summaries is on OSF. No quality appraisal score or double screening is reported.

\cf{Definition of a WAT} A philosophical approach with essential properties (worn on the body, sensors recording physiological/environmental data, electronic, forms vary, data analysis available to the user without a health professional, actionable information derived from sensors) and accidental ones (upload to server, docking, visualisation on app or device, immediate feedback) \pc{niksirat2024}{5}{To be a WAT, an object needs to have all of the following essential properties}. Smartwatches count; medical-connected devices (insulin pumps) and special-purpose devices do not. Ecosystem: WAT, companion app, provider servers, third-party applications (TPAs) such as Strava, connected devices.

\cf{Utility (Section 3, summary)} Benefits: motivation and increased physical activity (not always sustained), health improvement, medical benefits, social bonding, sleep awareness; usage patterns (power, intermittent, hop-on/hop-off, quantified-selfers); sharing and usability. Not needed in detail for the thesis.

\cf{Inferences from WAT data (Section 4)} ML can infer health metrics (ECG, post-surgery complications, SARS-CoV-2 infection, mental-health states such as stress resilience, cancer readmission), activities (eating, drinking, smoking, drunkenness), personal characteristics (height, weight, handwriting, mood, three of the Big Five traits), and location (inference of more than four-fifths of locations despite endpoint privacy zones). ``A large majority of research on WAT inferences show that WAT data can also be used in adversarial settings'', by insurers, employers and advertisers, to target or discriminate \pc{niksirat2024}{12}{A large majority of research on WAT inferences show that WAT data can also be used in adversarial settings}.

\cf{User awareness, concern, attitude, behaviour} \emph{Awareness}: ``Most users tend to have a limited understanding of how WATs and their ecosystem work'' and do not know who has access, how data are transmitted, stored and used; mental models do not match the TPA ecosystem; only a small share understood the data path to the provider \pc{niksirat2024}{12}{Most users tend to have a limited understanding of how WATs and their ecosystem work}; users believe sensitive information cannot be inferred from sensors they cannot see, and that most privacy risks are unlikely to materialise \pc{niksirat2024}{13}{Users tend to believe that most privacy risks are unlikely to materialize}. \emph{Concern}: most express only minor concerns \pc{niksirat2024}{13}{Most users express only minor privacy concerns}; concerns grow after misuse or when confronted with private-life questions (the privacy paradox); three concern types: data collection and storage, control over data (secondary purposes, sharing, forced-choice dilemma, post-purchase lock-in when policies change), storage security \pc{niksirat2024}{14}{Three types of concerns are recognized}. \emph{Attitudes}: location is the most sensitive data type; weight, sleep and PII/financial data also; willingness to share is inversely proportional to recipient-group size \pc{niksirat2024}{14}{Schneegass et al. [214] found that users’ willingness to share is inversely proportional to the size of the recipient group they share the data with}; families (monitoring children can erode trust), and intimate-partner abuse via WATs \pc{niksirat2024}{14}{Leitão [136] shows that WATs can be used by abusive partners for stalking, threatening, and harassing}. \emph{Behaviour}: few adjust privacy settings; almost one-third do not revoke TPA access because they forget it \pc{niksirat2024}{15}{Almost one-third of users usually do not revoke the WAT data access they granted to TPAs, as they forget that they installed them on their devices}; coping after violation is emotion- or problem-focused depending on threat and efficacy; individual differences (age, gender, region inconsistent); user classes (non-sensitive/sensitive; unconcerned/somewhat/highly concerned; users/former/non-users).

\cf{Policies, regulation, ethics} Policy issues: lack of legal accountability; usability (lengthy, complex; acceptance is a binary forced choice) \pc{niksirat2024}{16}{Privacy policies are lengthy, complex, and annoyingly profuse, thus users often do not read them to avoid cognitive load}; Hutton et al.\ found WAT apps met fewer privacy heuristics than other self-tracking apps \pc{niksirat2024}{16}{apps related to WATs generally met fewer heuristics}; proposed fixes: shorter policies, visualisation (Poli-see), partial consent (CURE), privacy notifications. Law: in the US WAT data are not protected health information because providers are not covered entities \pc{niksirat2024}{16}{WAT data is not counted as protected health information because SPs are not covered entities, unlike hospitals or clinics}; FDA treats WATs as low-risk wellness products; in the EU the GDPR gives better protection and applies to foreign manufacturers \pc{niksirat2024}{17}{the GDPR (General Data Protection Regulation) provides better protection for users}. Forensic use of WAT data (evidence; integrity and accuracy issues). Ethics: dataveillance and function creep \pc{niksirat2024}{18}{Lupton [150] uses the term dataveillance (i.e., digital surveillance of individuals) to explain how WAT use can lead to “function creep”}; insurer incentive programmes; research-data sharing. Workplace: wellness model versus performance-management model \pc{niksirat2024}{18}{Companies intending to adopt WAT-based wellness programs follow either a wellness model or a performance management model}, social pressure and boundary erosion.

\cf{PETs and joint analysis} Users are generally open to PETs \pc{niksirat2024}{19}{Users, in general, are open to using PETs}: anonymisation (accelerometer data can be de-anonymised), limited sharing/data minimisation, pedagogical tools (a video explainer improved attitudes while text did not), personalised setting recommenders, access-management with purpose/visibility/granularity/retention. Privacy-calculus: users weigh benefits and risks \pc{niksirat2024}{20}{technology users always weigh the perceived benefits and (privacy) risks}; utility often outweighs concern, but acceptable inference depends on context and country norms, e.g.\ religion and sexual orientation \pc{niksirat2024}{20}{Velykoivanenko et al. [242] argue that users’ concerns about the inference of certain types of information (e.g., religion and sexual orientation) are heavily dependent on the social norms and conditions in their country of residence}. Section 6 covers attacks (side channels, protocols, countermeasures) and Section 7 six research gaps; not summarised here.

\cf{Limits (card)} The review scope is WATs, so extrapolation to mental-health or cycle apps is not warranted; one author screened, no double screening or agreement, no risk-of-bias appraisal, no PRISMA flow numbers beyond 688/236 in the extract; the authors themselves note some findings are about 2012-era devices and policies that may have changed; the narrative synthesis does not weight studies by sample or design; some claims rest on single studies and are repeated as ``users believe''.

\cf{Use for this thesis} (1) A source for candidate inference-based harms (health, mental state, activity, location, personality) and for recipients that exploit them (insurers, employers, advertisers, abusive partners), to be linked to the harm tree with the caveat that evidence comes from wearables. (2) User-perception evidence for the harm stage: limited awareness of data flows and TPAs, forgotten third-party access, forced-choice consent, and lock-in after policy changes are plausible things to check in app policies and consent flows. (3) A model of a documented review (search string, sources, counts, eligibility, who screened, synthesis guideline) to compare with the thesis' own selection reporting; the thesis' stage-D protocol should add double screening or document why not. (4) Pairs with Gabriele and Chiasson (survey of tracker users), Hutton et al.\ (policy heuristics) and Cao et al.\ (vignettes).
\end{papercard}
